\subsection{Consequences of a proved twist exponent}
\label{subsec:actual-exponent}

Suppose the actual twists of $Z(\mathcal C)$ have order dividing a
proved integer $N$. In this subsection $\mathcal C$ is pseudo-unitary,
$D=\operatorname{FPdim}(\mathcal C)$, and all indicators are actual
indicators, before relative trace projection. The center formula makes
$\nu_k(X)$ periodic in $k$ modulo $N$. It also gives
\begin{equation}\label{eq:actual-indicator-bound}
 |\nu_k(X)|\le D^{-1}\sum_Z[F(Z):X]d_Z=d_X.
\end{equation}
The dimension identity on the right is
$\operatorname{FPdim}I(X)=Dd_X$.

\begin{proposition}[Full rotation spectra and prime recovery]
\label{prop:actual-rotation-prime}
Let $h_n(X)=[X^{\otimes n}:\mathbf1]$.
\begin{enumerate}
\item If $n\mid N$, the entire character of rotation on
$\operatorname{Hom}(\mathbf1,X^{\otimes n})$ is determined by the actual
center variables. Its value at the identity is $h_n(X)$. For
$0<r<n$, put $d=\gcd(n,r)$ and $m=n/d$. Then
\begin{equation}\label{eq:full-rotation-character}
 t_r=\operatorname{Tr}((E_X^{(n)})^r)
     =\sigma_{r/d}\!\left(\nu_m(X^{\otimes d})\right),
 \qquad
 c_e=\frac1n\sum_{r=0}^{n-1}t_r\zeta_n^{-er}
       \in\mathbb Z_{\ge0}\quad(0\le e<n).
\end{equation}
Here $\sigma_{r/d}$ acts on $\mathbb Q(\zeta_m)$; any unit lift modulo
$N$ gives the same answer.
\item If $p\nmid N$ is prime, then
\begin{equation}\label{eq:prime-indicator-congruence}
 \nu_p(X)\in\mathbb Z,\qquad
 \nu_p(X)\equiv h_p(X)\pmod p,\qquad |\nu_p(X)|\le d_X.
\end{equation}
Consequently, for an integer bound $B_X\ge d_X$ and a prime
$p>2B_X$, the residue $h_p(X)\bmod p$ allows at most one value of
$\nu_p(X)$ in $[-B_X,B_X]$. If no such value exists, categorification
is impossible. Values obtained from primes in the same residue class
modulo $N$ must agree.
\end{enumerate}
\end{proposition}
\begin{proof}
Lemma~\ref{lem:blockrotation} identifies $(E_X^{(n)})^d$ with rotation
on $m$ blocks of length $d$. This operator has $m$th roots as eigenvalues.
Raising it to the power $r/d$, which is coprime to $m$, applies the
cyclotomic automorphism $\sigma_{r/d}$ to its trace. This proves the
first formula. The reduction map from units modulo $N$ to units modulo
$m\mid N$ is surjective, and its kernel fixes $\mathbb Q(\zeta_m)$,
so the choice of lift is immaterial. Fourier inversion gives the actual
eigenvalue multiplicities $c_e$, proving their integral positivity.

For a prime $p\nmid N$, the indicator is both a rotation trace in
$\mathbb Q(\zeta_p)$ and a center trace in $\mathbb Q(\zeta_N)$.
The center field statement follows from its canonical congruence
representation \cite[Theorem~II]{DLN}. These fields intersect in
$\mathbb Q$, so it is a rational algebraic integer. A rational character of $C_p$ has equal multiplicity $b$ at
each primitive $p$th root. If the multiplicity at one is $a$, then
$h_p=a+(p-1)b$ and $\nu_p=a-b=h_p-pb$. This proves the congruence;
\eqref{eq:actual-indicator-bound} proves the bound. An interval of
length $2B_X<p$ contains at most one integer in a specified residue
class. Finally, periodicity modulo $N$ proves the compatibility between
primes. All $h_p\bmod p$ are obtained from based-ring matrix powers
modulo $p$; a maximum row sum of a fusion matrix supplies an integer
Perron--Frobenius bound $B_X$. No indicators are guessed.
\end{proof}

The same argument shows why an averaged primitive-root bucket cannot
always be treated as an arbitrary integer count: for $p\nmid N$ its
total multiplicity is divisible by $p-1$. Nonnegative real relaxations
lose this divisibility, but exact modular row certificates retain it.

\begin{proposition}[Coprime higher Gauss sums of the center]
\label{prop:center-higher-gauss}
Write $\tau_n=\sum_Z d_Z^2\theta_Z^n$. For every $n$ coprime to $N$,
\begin{equation}\label{eq:center-higher-gauss}
 \tau_n=\frac{D^2}{\sigma_{n^{-1}}(D)},
\end{equation}
where the inverse is modulo $N$.
\end{proposition}
\begin{proof}
Use the canonical anomaly-one representation of the center, so that
$\tau_1=D$, $S_{0Z}=d_Z/D$, and $T_{ZZ}=\theta_Z$.
The signed Galois permutation $\pi_n$ satisfies
\[
 \sigma_n(d_Z)^2=\frac{\sigma_n(D)^2}{D^2}d_{\pi_n(Z)}^2,
 \qquad \theta_{\pi_n(Z)}=\theta_Z^{n^2}
\]
by the Galois formulas of \cite[Theorem~II]{DLN}. Hence
\[
 \sigma_n(\tau_n)
 =\frac{\sigma_n(D)^2}{D^2}
     \sum_Zd_{\pi_n(Z)}^2\theta_{\pi_n(Z)}
 =\frac{\sigma_n(D)^2}{D}.
\]
Applying $\sigma_{n^{-1}}$ gives
\eqref{eq:center-higher-gauss}. These are exact linear equations in the
actual restriction-and-twist counts. They do not require an independent
choice of higher central charges.
\end{proof}
